\begin{lemma}[The fixed-time testing identity for the centered Ornstein--Uhlenbeck evolution]
For $s > 0$ write $r_s = e^{-s}$ and $\tau_s = \sqrt{1-e^{-2s}}$, and for
$g : \mathbb{R}^n \to \mathbb{R}$ write
$(A_s g)(x) = \int g(r_sx+\tau_sz)\,d\gamma_n(z)$ for the Mehler average of
\noderef{MehlerFormula}.  Let $\lambda$ denote Lebesgue measure on
$\mathbb{R}^n$.  Let $f : \mathbb{R}^n \to \mathbb{R}$ be smooth with compact
support and square $\gamma_n$-integrable, and write $u \in L^2(\gamma_n)$ for the
class it determines.  Let $\Pi_0 = L^0$ be the mean-zero projection of
\noderef{FracPow} and let $P_s$ be the Ornstein--Uhlenbeck semigroup of
\noderef{OUSemigroup}.  Let $\varphi : \mathbb{R}^n \to \mathbb{R}$ be smooth
with compact support, let $i \in \{1,\ldots,n\}$ and let $s > 0$.  Then:
\begin{enumerate}
\item the function
$x \mapsto (P_s\Pi_0u)(x)\,\partial_i\varphi(x)$ is $\lambda$-integrable;
\item the function
$x \mapsto \bigl(A_s(\partial_if)\bigr)(x)\,\varphi(x)$ is $\lambda$-integrable;
and
\item
\[
  \int_{\mathbb{R}^n}(P_s\Pi_0u)(x)\,\partial_i\varphi(x)\,d\lambda(x)
    = -e^{-s}\int_{\mathbb{R}^n}
        \bigl(A_s(\partial_if)\bigr)(x)\,\varphi(x)\,d\lambda(x) .
\]
\end{enumerate}

Here $(P_s\Pi_0u)(x)$ means the value at $x$ of the function to which the
$L^2(\gamma_n)$ class $P_s\Pi_0u$ is coerced.  Such a coercion is determined only
up to a $\gamma_n$-null set, but $\gamma_n$ and $\lambda$ have the same null sets
on $\mathbb{R}^n$ --- this is recorded as a separate fact below and used in the
proof --- so neither the integrability in clause (1) nor the value of the
integral in clause (3) depends on which coercion is taken.

This is the reference paper's fixed-time half of the derivation of the semigroup
formula for $R_\gamma$: the paper differentiates Mehler's formula to get
$DP_sf = e^{-s}P_sDf$ at each fixed $s$, and only afterwards passes to the time
integral.  The identity above is that fixed-time statement in tested form, with
the derivative moved onto the test function $\varphi$ by integration by parts and
with the mean-zero projection already accounted for; the factor $e^{-s}$ is
exactly the one the paper's fixed-time identity carries.  No time integral, no
subordination weight and no Riesz transform occurs here, and nothing is asserted
about the $s$-dependence of either side.
\end{lemma}
\begin{proof}
Throughout, $\lambda$ is Lebesgue measure and all integrals written
$\int_{\mathbb{R}^n}\cdots d\lambda$ are against $\lambda$; the Mehler averages
are against $\gamma_n$ in the inner variable.  Since $s > 0$ we have
$e^{-2s} < 1$, so $1-e^{-2s} > 0$ and therefore
\begin{equation*}
  \tau_s = \sqrt{1-e^{-2s}} > 0, \qquad\text{in particular } \tau_s \neq 0 .
  \tag{0}
\end{equation*}

\emph{Step 1: $\gamma_n$ and $\lambda$ have the same null sets.}
By clause (4) of \noderef{GaussLebesgueBallComparison}, for every
$A \subseteq \mathbb{R}^n$ one has $\gamma_n(A) = 0$ if and only if
$\lambda(A) = 0$.  Consequently a property holding $\gamma_n$-almost everywhere
holds $\lambda$-almost everywhere and conversely: the exceptional set is the same
set in both cases, and it is null for one measure exactly when it is null for the
other.  We use this in both directions below without further comment.

\emph{Step 2: regularity of the two Mehler averages.}
Apply \noderef{SemigroupGradientCommute} at this $s$, to this $f$ --- smooth with
compact support --- and at the index $i$.  Its clause (3) says that
$A_s(\partial_if)$ is continuous on $\mathbb{R}^n$.  Its clause (4) says that
$A_sf$ is differentiable at every point $x$ and that
\begin{equation*}
  \partial_i(A_sf)(x) = e^{-s}\bigl(A_s(\partial_if)\bigr)(x)
  \qquad\text{for every } x . \tag{1}
\end{equation*}
Being differentiable at every point, $A_sf$ is in particular continuous.  By
$(1)$ the function $x \mapsto \partial_i(A_sf)(x)$ is the constant $e^{-s}$ times
the continuous function $A_s(\partial_if)$, hence is itself continuous.

Also, $\varphi$ being smooth, $\partial_i\varphi$ is continuous; and
$\partial_i\varphi$ vanishes off the compact set $\operatorname{supp}\varphi$,
because a point outside that set has an open neighbourhood on which $\varphi$
vanishes identically and hence on which all derivatives of $\varphi$ vanish.  So
$\partial_i\varphi$ is continuous with compact support.

\emph{Step 3: the three integrability statements.}
A continuous function with compact support on $\mathbb{R}^n$ is
$\lambda$-integrable (shared context).  The products
\[
  (A_sf)\cdot\partial_i\varphi,
  \qquad
  c\cdot\partial_i\varphi \quad (c \in \mathbb{R} \text{ constant}),
  \qquad
  \bigl(A_s(\partial_if)\bigr)\cdot\varphi
\]
are each continuous, by Step 2, and each has compact support --- the first two
because $\partial_i\varphi$ does, the third because $\varphi$ does --- so all
three are $\lambda$-integrable.  The third is clause (2).  Being integrable, the
difference $\bigl((A_sf) - c\bigr)\cdot\partial_i\varphi$ is integrable as well,
for any constant $c$.

\emph{Step 4: the Mehler average of the class is the Mehler average of $f$.}
Let $N = \{\,y \in \mathbb{R}^n : u(y) \neq f(y)\,\}$, where $u$ is read through
its coercion.  By construction of the class of $f$ we have $u = f$
$\gamma_n$-almost everywhere, i.e.\ $\gamma_n(N) = 0$, so $\lambda(N) = 0$ by
Step 1.  Fix $x \in \mathbb{R}^n$ and put $\Psi(z) = r_sx+\tau_sz$.  Then
\[
  \Psi^{-1}(N)
    = \{\,z : r_sx+\tau_sz \in N\,\}
    = \bigl(z \mapsto \tau_sz\bigr)^{-1}
        \bigl(\{\,y : r_sx+y \in N\,\}\bigr) ,
\]
the two descriptions being the same set by definition of preimage.  Lebesgue
measure is invariant under translation (shared context), so
$\lambda(\{y : r_sx+y \in N\}) = \lambda(N) = 0$.  By $(0)$ we have
$\tau_s \neq 0$, so clause (2) of \noderef{LebesgueDilation}, applied with the
scalar $c = \tau_s$ and the $\lambda$-null set $\{y : r_sx+y \in N\}$, gives
$\lambda(\Psi^{-1}(N)) = 0$, hence $\gamma_n(\Psi^{-1}(N)) = 0$ by Step 1.
Outside $\Psi^{-1}(N)$ the two integrands $z \mapsto u(\Psi(z))$ and
$z \mapsto f(\Psi(z))$ agree, so they agree $\gamma_n$-almost everywhere and
their $\gamma_n$-integrals coincide:
\begin{equation*}
  \int u(r_sx+\tau_sz)\,d\gamma_n(z) = (A_sf)(x)
  \qquad\text{for every } x . \tag{2}
\end{equation*}
By clause (2) of \noderef{MehlerFormula}, applied at this $s$ to $v = u$, we have
$(P_su)(x) = \int u(r_sx+\tau_sz)\,d\gamma_n(z)$ for $\gamma_n$-almost every $x$.
Combining with $(2)$,
\begin{equation*}
  (P_su)(x) = (A_sf)(x)
  \qquad\text{for } \gamma_n\text{-a.e. } x,\text{ hence for }
  \lambda\text{-a.e. } x. \tag{3}
\end{equation*}

\emph{Step 5: the projection costs only a constant.}
Let $c = u_0$ be the Hermite coefficient of $u$ at the multi-index $\alpha = 0$,
relative to the basis of \noderef{HermiteBasis}; recall from that node that the
coefficient map $v \mapsto (v_\alpha)_\alpha$ is an isometric isomorphism onto
$\ell^2(\mathbb{N}_0^n)$, in particular injective.  By clause (1) of
\noderef{FracPowLinear} at exponent $0$, every element of $L^2(\gamma_n)$ lies in
$\operatorname{Dom}(L^0)$, so $\Pi_0u$ is defined; and by \noderef{FracPow} at
exponent $0$ its coefficients are
\[
  (\Pi_0u)_\alpha = \begin{cases} 0, & \alpha = 0,\\ u_\alpha, & \alpha \neq 0,\end{cases}
\]
the nonzero case using $|\alpha|^0 = 1$.  Write $h_0$ for the basis vector at
$\alpha = 0$; by clause (2) of \noderef{HermiteBasisVectorFacts} its coefficients
are $(h_0)_\alpha = 1$ at $\alpha = 0$ and $0$ at $\alpha \neq 0$, after
identifying $h_0$ with the corresponding Hermite element by clause (1) of that
node.  Now apply \noderef{OUSemigroupSpectral}, valid for every $s \geq 0$, to
$\Pi_0u$ and to $u$: the coefficients of $P_s\Pi_0u$ are
$e^{-s|\alpha|}(\Pi_0u)_\alpha$, and those of $P_su - c\,h_0$ are
$e^{-s|\alpha|}u_\alpha - c\,(h_0)_\alpha$.  At $\alpha = 0$ we have $|0| = 0$,
so the first is $e^{0}\cdot 0 = 0$ and the second is $e^{0}u_0 - c = 0$; at
$\alpha \neq 0$ the first is $e^{-s|\alpha|}u_\alpha$ and the second is
$e^{-s|\alpha|}u_\alpha - 0$.  The two coefficient families agree, so by
injectivity of the coefficient map
\begin{equation*}
  P_s\Pi_0u = P_su - c\,h_0 . \tag{4}
\end{equation*}
By \noderef{HermiteL2ZeroAEOne} the coercion of $h_0$ agrees $\gamma_n$-almost
everywhere with the constant function $1$ --- again using clause (1) of
\noderef{HermiteBasisVectorFacts} to identify $h_0$ with that Hermite element ---
so taking coercions in $(4)$ and using $(3)$,
\begin{equation*}
  (P_s\Pi_0u)(x) = (A_sf)(x) - c
  \qquad\text{for } \lambda\text{-a.e. } x. \tag{5}
\end{equation*}
By Step 3 the right-hand side of $(5)$, multiplied by $\partial_i\varphi$, is
$\lambda$-integrable; since two functions agreeing $\lambda$-almost everywhere
are integrable together, clause (1) follows from $(5)$.

\emph{Step 6: the constant term integrates to zero.}
We claim $\int_{\mathbb{R}^n}\partial_i\varphi\,d\lambda = 0$.  Apply
\noderef{CompIntegrationByParts} at the index $j := i$ with the following data.
Take $q = \mathrm{id}_{\mathbb{R}}$ and $q' \equiv 1$: then $q$ is differentiable
at every $t$ with derivative $q'(t) = 1$, and $q'$, being constant, is
continuous.  Take $w \equiv 1$, the constant function on $\mathbb{R}^n$: a
constant function is differentiable at every point with vanishing derivative, so
$\partial_i w(x) = 0$ for every $x$ and $x \mapsto \partial_i w(x)$ is the zero
function, which is continuous.  Take the test function to be $\varphi$, which is
smooth with compact support.  Both hypotheses on $w$ hold, and the conclusion of
that node reads
\[
  \int_{\mathbb{R}^n}q(w(x))\,\partial_i\varphi(x)\,d\lambda(x)
    = -\int_{\mathbb{R}^n}
        \bigl(q'(w(x))\,\partial_iw(x)\bigr)\varphi(x)\,d\lambda(x) .
\]
Here $q(w(x)) = 1$ and $q'(w(x))\,\partial_iw(x) = 1\cdot 0 = 0$ for every $x$,
so the left side is $\int\partial_i\varphi\,d\lambda$ and the right side is
$-\int 0\,d\lambda = 0$.  Hence
\begin{equation*}
  \int_{\mathbb{R}^n}c\,\partial_i\varphi\,d\lambda
    = c\int_{\mathbb{R}^n}\partial_i\varphi\,d\lambda = 0 . \tag{6}
\end{equation*}

\emph{Step 7: integration by parts at fixed time.}
Apply \noderef{CompIntegrationByParts} again at the index $j := i$, now with
$w = A_sf$, with $q = \mathrm{id}_{\mathbb{R}}$ and $q' \equiv 1$ as in Step 6,
and with the test function $\varphi$.  The two hypotheses that node places on $w$
are differentiability at every point and continuity of the single partial
derivative $\partial_iw$; both were verified in Step 2.  Since
$q(w(x)) = (A_sf)(x)$ and
$q'(w(x))\,\partial_iw(x) = \partial_i(A_sf)(x)$ for every $x$, its conclusion
reads
\[
  \int_{\mathbb{R}^n}(A_sf)\,\partial_i\varphi\,d\lambda
    = -\int_{\mathbb{R}^n}\partial_i(A_sf)\,\varphi\,d\lambda
    = -e^{-s}\int_{\mathbb{R}^n}\bigl(A_s(\partial_if)\bigr)\varphi\,d\lambda ,
\]
the last equality by $(1)$ together with linearity of the integral, legitimate
because $\bigl(A_s(\partial_if)\bigr)\varphi$ is integrable by Step 3.  There is
no boundary term because $\varphi$ is compactly supported.

\emph{Step 8: conclusion.}
By $(5)$ the integrands $(P_s\Pi_0u)\,\partial_i\varphi$ and
$\bigl((A_sf) - c\bigr)\partial_i\varphi$ agree $\lambda$-almost everywhere, so
their integrals coincide.  Both $(A_sf)\,\partial_i\varphi$ and
$c\,\partial_i\varphi$ are integrable by Step 3, so the integral of their
difference splits, and
\[
  \int_{\mathbb{R}^n}(P_s\Pi_0u)\,\partial_i\varphi\,d\lambda
    = \int_{\mathbb{R}^n}(A_sf)\,\partial_i\varphi\,d\lambda
      - \int_{\mathbb{R}^n}c\,\partial_i\varphi\,d\lambda
    = -e^{-s}\int_{\mathbb{R}^n}\bigl(A_s(\partial_if)\bigr)\varphi\,d\lambda ,
\]
using Step 7 for the first term and $(6)$ for the second.  This is clause (3).
\end{proof}
